\titulopreliminar{RESUMEN}
\vspace{10.75pt}

En la actualidad, el flujo de trabajo de los tatuadores presenta una alta fragmentación tecnológica, lo que hace hace que los profesionales dividan sus tareas administativas y artísticas en múltiples sistemas no integados. La falta de herramientas unificadas genera ineficiencias a la hora de trabajar, ya sea adecuando el tatuaje en una zona establecida del cuerpo, en pérdidas de tiempo en el cálculo de presupuestos y problemas gestionando las citas con los clientes.

El proyecto propone el diseño y construcción de una plataforma web integral que unifica el flujo de trabajo del estudio de tatuajes. El sistema incorpora módulos de almacenamiento en la nube para bocetos, previsualización 3D sobre modelos anatómicos, gestión de inventario con cotizaciones automatizadas, y un agendador de citas integrando una API de mensajería para los recordatorios.

Desde un análisis crítico, la solución posee factibilidad técnica, operativa y económica. Al utilizar tecnologías web y librerías de código abierto, el sistema elimina la dependencia de licencias costosas y no requiere hardware especializado. También teniendo en cuenta que se cuenta con la ayuda de un tatuador experto que ayuda en el proceso de captura de requerimientos y pruebas del sistema en cada iteración a través de su retroalimentación.

Como proyección de impacto, es un producto mínimo viable (MVP) que optimiza el modelo de negocio, con un objetivo de disminución de ausentismo de clientes, reducción de tiempo administrativo en cálculos de costos e inventario y una disminución de tiempo invertido en pruebas de diseño sobre la piel.
